\chapter{Ausgänge, Optionen und Genauigkeit}
\label{kap:format-ausgaben}

\section{Ausgänge}
\label{sec:ausgaenge}

\code{"outputs"} nennt die Größen, die gespeichert werden:

\begin{json}
"outputs": ["v(in)", "v(out)", "v(c,e)", "i(V1)", "T - 273.15", "v(out)*i(V1)", "T[*]"]
\end{json}

\begin{itemize}
  \item \code{v(knoten)}, \code{v(a,b)}, \code{i(bauteil)}, Variablen und
    beliebige Ausdrücke daraus. Die AC-Analyse akzeptiert nur die
    einfachen Formen.
  \item Zweigströme gibt es für Spannungsquellen, Spulen und die
    Spannungsquellen \code{E}, \code{H} und \code{B} (mit \code{v}). Für
    andere Ströme eine 0-V-Quelle in Reihe schalten – oder bei einem
    Widerstand \code{v(a,b)/R} ausgeben.
  \item Ohne \code{outputs} werden alle Knotenspannungen, Zweigströme und
    Variablen der Hauptschaltung ausgegeben; das Innere von Instanzen nur,
    wenn es genannt ist (\code{v(X1.mid)}).
  \item Jede Analyse kann mit eigenem \code{"outputs"} die globale Liste
    ersetzen.
\end{itemize}

\section{Ergebnisdateien}

Für die Eingabedatei \datei{name.json} und den Ausgabeordner \datei{out}
entstehen:

\begin{center}
\begin{tabularx}{\textwidth}{@{}lL@{}}
  \toprule
  Datei & Inhalt \\
  \midrule
  \datei{out/name.dc.csv} & DC-Sweep: Wert der Quelle, dann die Ausgänge \\
  \datei{out/name.ac.csv} & \code{freq}, dann je Ausgang \code{mag(...)} und \code{phase(...)} \\
  \datei{out/name.tran.csv} & \code{time}, dann die Ausgänge \\
  \datei{out/name.noise.csv} & Frequenz, Rauschdichten, Verstärkung \\
  \datei{out/name.pss.csv} & die eingeschwungene(n) Periode(n) \\
  \datei{out/name.tran.fourier.csv} & Fourier-Analyse: Harmonische mit Betrag und Phase \\
  \datei{out/name.<a>.measure.csv} & Messgrößen je Lauf einer Studie \\
  \datei{out/name.sens.csv} & Tabelle der Empfindlichkeiten \\
  \datei{out/name.fit.csv} & Messdaten und angepasste Modellkurve \\
  \datei{out/name.op.csv} & Arbeitspunkt einer Studie, eine Zeile je Lauf \\
  \bottomrule
\end{tabularx}
\end{center}

Statt des Analysetyps steht der \code{"name"} der Analyse im Dateinamen,
wenn sie einen hat. Alle Dateien sind CSV mit Komma als Trenner, Punkt als
Dezimalzeichen und einer Kopfzeile:

\begin{shell}
time,v(in),v(out)
0,0,0
4e-06,0,0
...
\end{shell}

Der einfache Arbeitspunkt steht nur auf der Konsole (und in der
Zusammenfassung, Abschnitt~\ref{sec:summary}).

\section{Optionen}
\label{sec:optionen}

\begin{json}
"options": {
  "reltol": 1e-4,       // relative Toleranz: Newton und Fehler je Zeitschritt
  "vntol": 1e-6,        // absolute Toleranz für Spannungen [V]
  "abstol": 1e-12,      // absolute Toleranz für Ströme [A]
  "vartol": 1e-9,       // absolute Toleranz für Variablen
  "lte_floor": 0.1,     // Bezugsgröße des Zeitschrittfehlers
  "trtol": 1,           // Faktor auf den zulässigen Zeitschrittfehler
  "gmin": 1e-12,        // Leitwert parallel zu jedem pn-Übergang [S]
  "temp": 27,           // Temperatur der Schaltung [°C]
  "tnom": 27,           // Temperatur der Modellparameter [°C]
  "maxiter_dc": 100,    // Newton-Iterationen im Arbeitspunkt
  "maxiter_tran": 20,   // Newton-Iterationen je Zeitschritt
  "reuse_jacobian": true,   // zerlegte Matrix im Zeitschritt weiterverwenden
  "index_reduction": true   // Systeme mit Index über 2 umformen
},
"nodeset": { "out": 2.5 },   // Startwerte für den Arbeitspunkt
"ic": { "c1": 5 }            // Knotenspannungen bei t = 0 für "uic"
\end{json}

Die angegebenen Werte sind die Standardwerte.

\begin{description}[style=nextline,leftmargin=2em]
  \item[{\code{nodeset}}]
    Startwerte von Knotenspannungen für die Suche des Arbeitspunkts. Sie
    helfen bei schwer konvergierenden Schaltungen und wählen bei bistabilen
    Schaltungen (Flipflops, Schmitt-Trigger) den Zustand. Der Arbeitspunkt
    selbst ist dann eine echte Lösung; die Werte sind nur der Anfang.
  \item[{\code{ic}}]
    Knotenspannungen bei $t=0$, nur mit \code{"uic": true} in der
    Transientenanalyse.
  \item[{\code{gmin}}]
    Ein kleiner Leitwert parallel zu jedem pn-Übergang, der die Matrix auch
    bei gesperrten Übergängen gut lösbar hält.
  \item[{\code{reuse_jacobian}}]
    Innerhalb eines Zeitschritts wird die zerlegte Matrix der ersten
    Newton-Iteration weiterverwendet, solange jede Iteration die Korrektur
    mindestens halbiert. Das spart rund 20\,\% Rechenzeit; \code{false}
    schaltet es ab.
\end{description}

\section{Genauigkeit der Zeitintegration}
\label{sec:genauigkeit}

Pro Schritt darf der geschätzte Fehler einer Unbekannten
\[
  \mathit{reltol}\cdot\max\bigl(|x|,\ \mathit{lte\_floor}\cdot\text{Spitzenwert von }|x|\bigr)
  + \text{absolute Toleranz}
\]
betragen (multipliziert mit \code{trtol}).

\begin{itemize}
  \item \code{lte_floor: 0} misst den Fehler relativ zum momentanen Wert.
    Das ist die übliche Definition und richtig für Lösungen, die über
    Größenordnungen abklingen; bei schwingenden Signalen kostet jeder
    Nulldurchgang kleine Schritte.
  \item \code{lte_floor: 1} misst relativ zum bisherigen Spitzenwert: am
    schnellsten, für Signale mit gleichbleibender Amplitude ausreichend.
  \item Der Standard 0,1 liegt dazwischen.
\end{itemize}

Das ist eine Schranke für den Fehler \emph{pro Schritt}. Über viele Schritte
summiert er sich; für genauere Ergebnisse \code{reltol} verkleinern. Mit dem
Standardverfahren \code{bdf} kostet das wenig, weil es dann höhere Ordnungen
nutzt. Ein Beispiel: Der Kippschwinger in
\datei{examples/relaxation_oscillator.json} trifft seine Periode mit den
Standardwerten auf 0,004\,\%, mit \code{reltol} $=10^{-9}$ und \code{vntol}
$=10^{-10}$ auf 1\,ns.

\begin{hinweis}[Faustregeln]
  \begin{itemize}[leftmargin=1.2em]
    \item Kurven sehen eckig aus: \code{tstep} kleiner wählen (begrenzt die
      Schrittweite) oder \code{output_step} für ein feines Ausgaberaster.
    \item Ergebnis soll genauer werden: \code{reltol} um eine oder zwei
      Größenordnungen verkleinern.
    \item Schwingkreis hoher Güte verliert Amplitude: \code{"method":
      "trap"}.
    \item Rechnung zu langsam bei vielen Nulldurchgängen:
      \code{lte_floor} auf 1.
  \end{itemize}
\end{hinweis}
